%!TEX program = xelatex
% !TEX encoding = ISO-8859-1
\def\ano{2020}
\def\semestre{1}
\def\disciplina{\'Algebra 1}
\def\turma{C}
\def\autor{Jos\'e Ant\^onio O. Freitas}
\def\instituto{MAT-UnB}

\documentclass{beamer}
\usetheme{Madrid}
\usecolortheme{beaver}
% \mode<presentation>
\usepackage{caption}
\usepackage{textpos}
\usepackage{amssymb}
\usepackage{amsmath,amsfonts,amsthm,amstext}
\usepackage[brazil]{babel}
% \usepackage[latin1]{inputenc}
\usepackage{graphicx}
\graphicspath{{/home/jfreitas/GitHub_Repos/video_aulas/logo/}{C:/Users/josea/Documents/GitHub_Repos/video_aulas/logo/}}
\usepackage{enumitem}
\usepackage{multicol}
\usepackage{answers}
\usepackage{tikz,ifthen}
\usetikzlibrary{lindenmayersystems}
\usetikzlibrary[shadings]

\renewcommand{\qedsymbol}{$\blacksquare$}
\newtheorem{definicao}{Defini\c{c}\~ao}[section]
\newtheorem{definicoes}{Defini\c{c}\~oes}[section]
\newtheorem{exemplo}{Exemplo}[section]
\newtheorem{exemplos}{Exemplos}[section]
\newtheorem{exercicio}{Exerc{\'\i}cio}
\newtheorem{observacao}{Observa{\c c}{\~a}o:}[section]
\newtheorem{observacoes}{Observa{\c c}{\~o}es:}[section]
\newtheorem*{solucao}{Solu{\c c}{\~a}o:}
\newtheorem{proposicao}{Proposi\c{c}\~ao}
\newtheorem{corolario}{Corol\'ario}
\newenvironment{prova}[1][Prova]{\noindent\textbf{#1:} }{\qedsymbol}%{\ \rule{0.5em}{0.5em}}
\newcommand{\nsub}{\varsubsetneq}
\newcommand{\vaz}{\emptyset}
\newcommand{\im}{{\rm Im\,}}
\newcommand{\sub}{\subseteq}
\newcommand{\n}{\mathbb{N}}
\newcommand{\z}{\mathbb{Z}}
\newcommand{\rac}{\mathbb{Q}}
\newcommand{\real}{\mathbb{R}}
\newcommand{\complex}{\mathbb{C}}
\newcommand{\cp}[1]{\mathbb{#1}}
\newcommand{\ch}{\mbox{\textrm{car\,}}\nobreak}
\newcommand{\vesp}[1]{\vspace{ #1  cm}}
\newcommand{\compcent}[1]{\vcenter{\hbox{$#1\circ$}}}
\newcommand{\comp}{\mathbin{\mathchoice
{\compcent\scriptstyle}{\compcent\scriptstyle}
{\compcent\scriptscriptstyle}{\compcent\scriptscriptstyle}}}

\title{Rela\c{c}\~ao de Equival\^encia}
\author[\autor]{\autor}
\institute[\instituto]{\instituto}
\date{\today}

\begin{document}
    \begin{frame}
        \maketitle
    \end{frame}

    \logo{\includegraphics[scale=.1]{logo-MAT.png}\vspace*{8.5cm}}

    \begin{frame}
        \begin{definicao}
            Seja $A$ um conjunto n{\~a}o vazio \pause e $R\subseteq A \times A$. \pause Dizemos que $R$ \pause {\'e} uma \textbf{rela{\c c}{\~a}o de equival{\^e}ncia} se:\pause
            \begin{enumerate}[label={\roman*})]
                \item Para todo $x \in A$, \pause $(x,x) \in R$. \pause \textit{(Propriedade Reflexiva)}\pause
                \item Se $(x, y) \in R$, \pause ent\~ao $(y, x) \in R$. \pause \textit{(Propriedade Sim\'etrica)}\pause
                \item Se $(x, y) \in R$ \pause e $(y, z) \in R$, \pause ent\~ao $(x, z)\in R$. \pause \textit{(Propriedade Transitiva)}\pause
            \end{enumerate}
        \end{definicao}

        Quando $R\subseteq A \times A$ {\'e} uma rela{\c c}{\~a}o de equival{\^e}ncia, \pause dizemos que $R$ {\'e} uma rela{\c c}{\~a}o de equival{\^e}ncia em $A$. \pause Quando dois elementos $x$, $y \in A$ \pause s{\~a}o tais que $(x,y) \in R$, \pause dizemos que $x$ e $y$ \textbf{s{\~a}o relacionados} \pause ou que $x$ e $y$ \textbf{est\~ao relacionados}.\pause
    \end{frame}

    \begin{frame}
        \begin{exemplos}
            \begin{itemize}
                \item[1)] Seja A=\{1,2,3,4\}. \pause Temos
                \begin{center}
                    \begin{tabular}{l}
                        $A\times A = \{(1,1); (1,2);(1,3);(1,4);(2,1);(2,2);(2,3);(2,4);$\\
                        $(3,1);(3,2);(3,3);(3,4);(4,1);(4,2);(4,3);(4,4)\}.$
                    \end{tabular}
                \end{center}
                Quais dos seguintes conjuntos s\~ao exemplos de rela{\c c}{\~o}es de equival{\^e}ncia?
                \begin{itemize}
                    \item $R_0 = \emptyset$\pause \vspace{.3cm}
                    \item $R_{1}= A\times A$\pause \vspace{.3cm}
                    \item $R_{2}=\{(1,1);(2,2);(3,3)\}$\pause \vspace{.3cm}
                    \item $R_{3}=\{(1,1);(2,2);(3,3);(4,4);(1,2);(2,1)\}$\pause \vspace{.3cm}
                    \item $R_{4}=\{(1,1);(2,2);(3,3);(4,4)\}$\pause \vspace{.3cm}
                    \item $R_{5}=\{(1,1);(2,2);(3,3);(4,4);(1,2);(2,1);(2,4);(4;2)\}$\pause
                \end{itemize}
            \end{itemize}
        \end{exemplos}
    \end{frame}
    \begin{frame}
    \end{frame}
    \begin{frame}
        \begin{exemplos}
            \begin{itemize}[label={\arabic*})]                
                \item[2
                )] Seja $A = \z$ \pause e $R\subseteq \z\times \z$ \pause definida por
                \begin{center}
                    $R = \{(x,y)\in \z \times \z \mid \pause x - y = 2k, \pause \mbox{ para algum } k \in \z\}$.\pause
                \end{center}
                 Mostre que $R$ \'e uma rela{\c c}{\~a}o de equival{\^e}ncia sobre $\z$.\pause
            \end{itemize}
        \end{exemplos}
        \vspace{3cm}
    \end{frame}
    \begin{frame}
    \end{frame}
    \begin{frame}
        \begin{exemplos}
            \begin{itemize}[label={\arabic*})]                
                \item[3)] Seja $A = \z \times \z^*$, onde $\mathbb{Z}^* = \mathbb{Z} \setminus \{0\}$. \pause  Para $(a, b)$, $(c, d) \in A$, \pause considere a seguinte rela{\c c}{\~a}o\pause
                \begin{center}
                    $(a, b) S (c, d) \pause \mbox{ quando } ad = bc.$
                \end{center}
                Mostre que $S$ \'e uma rela{\c c}{\~a}o de equival{\^e}ncia sobre $A$.\pause
            \end{itemize}
        \end{exemplos}
        \vspace{3cm}
    \end{frame}
    \begin{frame}
    \end{frame}
\end{document}